\subsection{Adaptive Universal Samplers}
\label{sec:universal-samplers}

A universal sampler turns one public setup into a source of publicly
reproducible samples from many distributions.  To specify a distribution,
we give a circuit $d$ that maps random bits to a sample.  After setup has
produced public parameters $U$, anyone can compute
$\Sample(U,d)$.  Everyone using the same description obtains the
same sample, but the random tape is not exposed separately from the
sample.  Distinct descriptions should behave as if they had been evaluated
on independent random tapes.

We use the adaptive definition of Hofheinz et al.\
\cite[full version, Section~3.2, Definition~4]{HofheinzEtAl2016UniversalSamplers}.
Here, adaptive means that descriptions may be chosen after seeing $U$ and
the results of earlier computations.

\subsubsection{Syntax}
\label{sec:universal-sampler-syntax}

Let $\lambda$ be the security parameter; efficient algorithms run in
probabilistic polynomial time.  Fix polynomial bounds
$\ell=\ell(\lambda)$, $m=m(\lambda)$, and $k=k(\lambda)$ for the
description length, number of random bits, and sample length, respectively.
A description $d\in\bits^\ell$ specifies a deterministic circuit
\[
  d:\bits^m\longrightarrow\bits^k.
\]
Shorter descriptions and outputs are padded, and circuits using fewer
than $m$ random bits ignore the surplus.  We fix a total decoding
convention: any invalid $\ell$-bit description denotes the circuit that
always outputs $0^k$.  The exact description string identifies the
distribution request; two different descriptions may compute the same
function without identifying the same request.

A universal sampler consists of two efficient public oracle algorithms:
\begin{description}
  \item[$\Setup^{H}(1^\lambda)\to U$.]
    The randomized setup algorithm produces the public sampler parameters.
  \item[$\Sample^{H}(U,d)\to p_d$.]
    The sampling algorithm computes a $k$-bit sample.  For fixed $U$, $d$,
    and oracle $H$, this computation is deterministic.
\end{description}
The bounds $(\ell,m,k)$ are fixed for the scheme and suppressed from its
algorithm inputs.  The superscript $H$ denotes access to a public random
oracle, with domain $\bits^*$ and a scheme-specified polynomial output
length.  It is available to setup, public sampling, and the adversary.
The distribution circuit $d$ itself has no oracle access.

\subsubsection{The ideal source of samples}
\label{sec:ideal-samples}

The ideal source keeps a private table of random tapes, initially empty.
Its command is
\[
  \IdealSample(d):\qquad
  \begin{cases}
    \text{if $r_d$ is undefined, choose $r_d\sample\bits^m$ independently;}\\
    \text{return $d(r_d)$.}
  \end{cases}
\]
Only the sample is returned, with no separate disclosure of $r_d$.
Thus a repeated description
returns the same value, whereas a new description receives a fresh random
tape.  Equivalently, $\IdealSample(d)=d(F(d))$ for a private truly
random function $F:\bits^\ell\to\bits^m$.  This $F$ is not the public
random oracle $H$.

The security definition must reconcile this ideal source with public
evaluation: even in the ideal experiment, the adversary receives $U$ and
can run $\Sample$ itself.  It is therefore not enough to replace
an answer at a sampling interface.  A simulator must arrange that actual
public evaluations agree with the ideal samples.


\subsubsection{Adversary commands and admissibility}
\label{sec:sampler-commands}

The adversary $\mathcal A$ is an efficient interactive algorithm given
$(1^\lambda,U)$, with the following commands:
\begin{description}
  \item[$\QueryRO(x)$.]
    Request the public-oracle value at $x\in\bits^*$.
  \item[$\Declare(d)$.]
    Declare a description $d\in\bits^\ell$.  There is no reply.
    The adversary then runs $p\gets\Sample^{\QueryRO}(U,d)$ and
    records $\langle\DeclareOutput,d,p\rangle$ on an auxiliary tape.
  \item[$\Finish(b)$.]
    Terminate with $b\in\bits$.  Termination without an output bit is
    interpreted as $\Finish(0)$.
\end{description}
Records use a fixed unambiguous encoding.  The adversary is
\emph{admissible} if it follows this sampling procedure: every record
$\langle\DeclareOutput,d,p\rangle$ contains the result of its declared
evaluation, so $p=\Sample^{\QueryRO}(U,d)$.  Its other local computations
and oracle queries are unrestricted.

\subsubsection{Real and ideal experiments}
\label{sec:sampler-experiments}

\paragraph{Real experiment $\Real_{\mathcal A}(1^\lambda)$.}
\begin{enumerate}
  \item Initialize a random oracle $H$: its first query at any string
    receives an independent uniform answer of the prescribed length;
    repeated queries receive the same answer.
  \item Compute $U\sample\Setup^{H}(1^\lambda)$ and give
    $(1^\lambda,U)$ to $\mathcal A$.  Retain all oracle answers made
    during setup.
  \item Run $\mathcal A$, answering every $\QueryRO(x)$ with $H(x)$.
    A $\Declare(d)$ command receives no reply.
  \item On $\Finish(b)$, output $b$.
\end{enumerate}

\paragraph{Ideal experiment $\Ideal_{\mathcal A,\mathcal S}
(1^\lambda)$.}
The simulator consists of two efficient algorithms
$\mathcal S=(\SimSetup,\SimRO)$.  Both have access to the
same private $\IdealSample$ oracle.  The algorithm
$\SimRO$ is stateful; its state persists across requests.
\begin{enumerate}
  \item Initialize the private random-tape table for
    $\IdealSample$.
  \item Compute
    $(U,\SimState)\sample
    \SimSetup^{\IdealSample}(1^\lambda)$.
    Give $(1^\lambda,U)$ to $\mathcal A$ and retain $\SimState$ as
    the simulator's private state.
  \item Run $\mathcal A$.  On each $\QueryRO(x)$, compute
    $(a,\SimState)\sample
    \SimRO^{\IdealSample}(\SimState,x)$ and return $a$.
  \item After each $\Declare(d)$, monitor the auxiliary tape for the
    resulting $\langle\DeclareOutput,d,p\rangle$ record.  Obtain
    $p^*\gets\IdealSample(d)$.  If $p\ne p^*$, terminate with
    the distinguished result $\Inconsistent$.
  \item On $\Finish(b)$, output $b$.
\end{enumerate}
The adversary has no direct access to $\IdealSample$.  The
simulator receives only the $\QueryRO$ requests, not the sample
declarations, recorded results, or consistency checker's oracle calls.
There is no separately sampled public oracle in this experiment:
$\SimRO$ supplies its answers.

\begin{definition}[Adaptive universal sampling]
\label{def:adaptive-universal-sampler}
The pair $(\Setup,\Sample)$ is an adaptively secure
universal sampler if there exist efficient simulator algorithms
\[
  \mathcal S=(\SimSetup,\SimRO)
\]
such that, for every
efficient admissible adversary $\mathcal A$, there is a negligible
function $\varepsilon_{\mathcal A}$ satisfying
\begin{align*}
  \left|\Pr[\Real_{\mathcal A}(1^\lambda)=1]
    -\Pr[\Ideal_{\mathcal A,\mathcal S}(1^\lambda)=1]\right|
    &\leq\varepsilon_{\mathcal A}(\lambda),\\
  \Pr[\Ideal_{\mathcal A,\mathcal S}(1^\lambda)
      =\Inconsistent]
    &\leq\varepsilon_{\mathcal A}(\lambda).
\end{align*}
\end{definition}

The first condition hides the simulation.  The second ensures that an
honest public evaluation actually produces the corresponding ideal sample;
it is the source definition's \emph{honest-sample violation} condition.
The result $\Inconsistent$ is reserved for this failure, not for
an adversary choosing to stop its own execution.

There is no setup-time bound on the number of commands.  Each efficient
adversary makes only polynomially many of them, but the scheme is not
initialized with that polynomial.  The fixed bounds $(\ell,m,k)$ apply to
individual descriptions and samples, not to their total number.

\subsubsection{Existence}

\begin{theorem}[Hofheinz et al.]
\label{thm:adaptive-universal-sampler}
If indistinguishability obfuscation and one-way functions exist, then
adaptively secure universal samplers exist in the random-oracle model.
\end{theorem}
This is \cite[full version, Theorem~3]{HofheinzEtAl2016UniversalSamplers}.
We invoke this theorem as a black box and do not reproduce its
construction or proof.
